% Folder 57, batch 03 (archivists' pages 41-60).
% Pass: Opus 5.5 (claude-opus-5-5), 2026-10-03 — first pass, unchecked against the pages by a human.
%
% Grothendieck's own manuscript, paginated by him 4-11 then a new run « Schéma
% d'Albanese ». Fast hand: statements come through, connective prose is often lost.
\documentclass[11pt,a4paper]{article}
\input{../preamble/grothendieck}

\folder{57}
\batch{3}
\pages{41}{60}
\foldertitle{Picard : tapuscrit annoté (s.d.), lettres (s.d., 1962).}
\dating{1962-[vers 1968]}
\watermark{Édition de démonstration}

\begin{document}

\section{Théorème de rigidité}

\page{42}\note{p.~4 de l'auteur. Le début de l'exposé (pp.~1--3 de l'auteur) précède ce lot.}
\marginal{Déf.}
Un préschéma en groupes commutatif $G$ de type fini sur $S$ est dit
\emph{rigide} si \struck{pour tout}
\begin{enumerate}
  \item[(i)] $G$ est \struck{plat} \ill{} sur $S$ ;
  \item[(ii)] Pour tout $s \in S$, et \struck{tout} \add{nombre} \struck{entier}
  premier $\ell$ \struck{à la caractéristique} distinct de car $k(s)$,
  $\overline{G}_s$ est \struck{\ill{} plus petit sous-schéma fermé \ill{} de
  $\overline{G}_s$ \ill{}} \ill{} l'adhérence des ${}_{\ell^k}\overline{G}_s$
  [où $\overline{G}_s$ \ill{} \ill{} \ill{} $k(s)$] ;
  \item[(iii)] Pour tout \struck{\ill{}} ouvert $U$ de $S$, \ill{} \ill{} \ill{}
  des pts de $U$, le schéma \ill{}$(G|U)$ est \uncertain{fini} sur $S$.
\end{enumerate}
\note{dans un cadre, à gauche de (ii) : « ${}_{\ell^k}(\overline{G}_s)$ ».}

Remarquons qu'on peut prouver que il \struck{\ill{}}
\marginal{Il suffit de l'exiger \ill{} \ill{}}

\marginal{Prop.}
suffit pour \struck{\ill{}} (iii) que \struck{$G_s$ \ill{} \ill{}} \add{\ill{}}
\struck{dans} les fibres \ill{} $G_a$, $G_m$, Ab, $\mathbb{Z}/\ell$,
$\mathbb{Z}/p$ qui \ill{} \ill{} $G_{s_2} = G$, \ill{} \ill{} \ill{} $G_a$, \ill{}
$\mathbb{Z}/p$. \struck{\ill{}} \struck{En effet,} \struck{(la seule condition
est \ill{} suffisante} \struck{Pour la nécessité}, \ill{} \ill{} \ill{}
qu'une extension de groupes rigides \ill{} \ill{} est rigide ; puisque $G_m$,
Ab sont rigides (\ill{} \ill{} \uncertain{classiques}), $\mathbb{Z}/\ell$
\ill{} \ill{}. \struck{\ill{}} \ill{} \ill{} la \ill{} de structure des
groupes commutatifs, \ill{} \add{\ill{}} \ill{}

\marginal{en face des groupes cités : « $G_m$ \quad Ab \quad $\mathbb{Z}_\ell$ / $G_a$ \quad $\mathbb{Z}_p$ »}
\marginal{\ill{} \ill{} \ill{} \ill{} fid. plat \ill{}}

\page{43}
on \struck{\ill{} de \ill{} \ill{}} \ill{} la \ill{}, \ill{} \ill{} \ill{} de
structure, à prouver qu'une extension d'un schéma abélien $A$ par $G_a$
\ill{} \ill{} rigide [\ill{}, \ill{} \ill{} quotient d'un rigide est rigide].
\emph{On \ill{} \ill{} \ill{}} : une extension \emph{non triviale} de $A$ par
$G_a$, du moins en car.~$0$, est rigide.

\textbf{Proposition.} Soit $G$ schéma en groupes commutatif \add{loc.} de type
fini sur $S$, \emph{rigide}. Soit $n$ un entier premier aux car.\ des corps
résiduels de $S$. Alors, \struck{\ill{} \ill{} $G$} pour tout \ill{} $U$ de
$G$, $U$ est le plus petit sous-schéma fermé de $U$ qui majore les
${}_{n^k}G \cap U$.

Démonstration donnée plus haut.

\textbf{Proposition.} Soit $G$ un $S$-préschéma en gr.\ \struck{\ill{}}
\add{loc.} \struck{\ill{}} commutatif, \add{\ill{}} de type fini sur $S$,
(soit \struck{\ill{}} $n$ un entier \struck{\ill{}} premier aux car.\ des
corps résiduels de $S$. Alors \add{le morphisme $G \xrightarrow{n} G$ \ill{}}
\add{est \uncertain{surjectif}, en particulier} \struck{\ill{}} ${}_nG$ est
\struck{\ill{}} \uncertain{non ramifié} sur $S$ ; \struck{\ill{}} si $G$ est
plat sur $S$, ${}_nG$ est \uncertain{étale} sur $S$.

\page{44}\note{p.~5 de l'auteur.}
En effet, comme $G$ est plat sur $S$, on est ramené \struck{\ill{}} à la
première assertion \add{en} \struck{\ill{}} \struck{différentielle d'\ill{}}
\add{$S = \operatorname{Spec}(k)$}, et \ill{} \ill{} \ill{} \ill{} $k$ alg.\
clos. Alors $G = G_{\mathrm{red}} \times G_{\mathrm{inf}}$, où $G_{\mathrm{inf}}$
est \ill{} \struck{\ill{}} $G$ est \ill{} sur $k$, donc \struck{\ill{}}
\struck{\ill{} aux car.\ \ill{}} \ill{} que la \uncertain{notion}
\uncertain{différentielle} : on est ramené à prouver que la
\struck{\ill{}} \ill{} \ill{} \ill{} est \uncertain{partout} surjective, ce
qui est immédiat, car elle s'identifie à la multiplication par $n$ [\ill{}
\ill{} le foncteur $\operatorname{Lie}(G)$ est additif, et \ill{} les groupes
\ill{}\ commutatifs sur $k$].
\marginal{\ill{} ${}_nG = {}_nG_{\mathrm{red}}$ \ill{} \ill{} \ill{} \ill{} p.~\ill{} \ill{}}

\textbf{Lemme.} Soient $S$ schéma \struck{\ill{}} loc.\ noeth.\ connexe, $X$ un
\struck{schéma \ill{}} $S$-préschéma \struck{\ill{}} \add{\uncertain{non ramifié}},
$X_1$ et $X_2$ deux sous-schémas \struck{\ill{}} de $X$ étales \add{et finis}
sur $S$. \struck{\ill{}} \struck{\ill{} $X$, \ill{} \ill{} $X_1$ et $X_2$,
finis sur $S$} Soit $s \in S$ tel que $X_{1s} = X_{2s}$, alors $X_1 = X_2$.
\marginal{Introduisant $X_1 \times_X X_2 = X'$, fini \uncertain{étale} sur $S$, on
est ramené au cas où $X_1 \subset X_2$, \ill{} \ill{} des pts \ill{}
\ill{} \ill{} $X_{2s}$ \ill{} \ill{}}

Supposons d'abord $S = \operatorname{Spec}(A)$, $A$ \ill{} local. On est ramené
\ill{} \ill{} \ill{} $A$ est complet, donc \struck{\ill{} $X$ \ill{} fini sur $S$}
\add{[comme \ill{}]} $X$ \ill{} \ill{} \add{\ill{} $X_i$ \ill{}} \struck{des}
schémas \ill{} finis sur $S$ \struck{et \ill{} \ill{} \ill{} \ill{} fibres
\ill{} sur $S$}] \ill{} \ill{} \ill{} $X$ est fini sur $S$ [\ill{} \ill{}
\ill{} $X_i$ \ill{} \ill{} \ill{} \ill{} sans \ill{}]

\page{45}
\add{\ill{} \ill{} \ill{} supposer que $X$ est un sous-schéma fermé de $S$
[en faisant une extension \struck{\ill{}} plate de $S$].} Mais alors $X_1$ et
$X_2$ étant finis \add{et étales} sur $S$, et \ill{} \ill{}, sont identiques à
$X$, donc identiques. \struck{l'assertion est bien connue.}

Dans le cas général, on veut dire que $X_1$ et $X_2$ coïncident \add{\ill{}}
\ill{} un voisinage ouvert \struck{\ill{}} de $s$. Soit $U$ le plus grand
ouvert \ill{} de $S$ tel que $X_1|U = X_2|U$, il reste à prouver que $U$ est
\emph{fermé}. Or si $s' \in \overline{U}$, \ill{} \ill{} $s' \in
\overline{s''}$, $s'' \in U$, donc
\[
  X_{is'} = \overline{X_{is''}} \cap X_{s'}
\]
[car $X_{s'} \to S$ \ill{} \ill{}] donc $X_{1s'} = X_{2s'}$, donc \ill{} \ill{}
\ill{} \ill{} précède $s' \in U$, cqfd.

\struck{\ill{}}

\textbf{Théorème de rigidité.} Soient $S$ un préschéma \add{loc.} noeth.\
\emph{connexe}, $G$ un schéma en groupes \add{\ill{}} \struck{\ill{}} de type
fini \struck{\ill{}} sur $S$, $G_1$ et $G_2$ deux sous-schémas en groupes
fermés \emph{rigides}, \uncertain{soit} $s \in S$ tels que $G_{1s} = G_{2s}$.
\struck{\ill{} \ill{} $s \in S$.} \struck{Alors} Alors $G_1 = G_2$.

\page{46}\note{p.~6 de l'auteur.}
\textbf{Corollaire 1.} Soient $G$, $H$ deux \struck{\ill{}} schémas en groupes
comm.\ de type fini sur $S$ loc.\ noeth.\ connexe, $G$ \emph{rigide}. Soient
$u, v \colon G \to H$ deux hom.\ de $S$-groupes, et $s \in S$ tel que
$u_s, v_s \colon G_s \rightrightarrows H_s$ soient égaux. Alors $u = v$.

On applique le th.\ de rigidité au graphe.

\textbf{Corollaire 2.} Soient $S$ un schéma loc.\ noeth.,
$S' \to S$ un morphisme \emph{fid.\ plat} \add{et de type fini}
\struck{\ill{}}, avec $S'$ loc.\ noeth., $G$ un $S$-schéma en groupes
\ill{} de type fini sur $S$, $H'$ un sous-schéma en groupes fermé de
$G' = G \times_S S'$. On suppose
\begin{enumerate}
  \item[(i)] les fibres de $S' \to S$ sont \uncertain{géom.} connexes,
  \item[(ii)] $H'$ est rigide.
\end{enumerate}
Alors il existe un sous-schéma en groupes \add{fermé} $H$ de $G$
(nécessairement unique, et rigide) tel que $H' = H \times_S S'$.

On applique le théorème de la descente, pour laquelle il faut seulement
\note{la phrase se poursuit p.~47.}

\page{47}
montrer que les deux images inverses de $H'$ dans $G'' = G \times_S S''$
(où $S'' = S' \times_S S'$) sont les mêmes.
\note{le passage qui suit, de « On peut supposer » à « il suffit pour cela de »,
est encadré puis biffé de trois longs traits obliques ; il est transcrit tel
qu'il se lit.}
\struck{On peut supposer $S$ \ill{}} \struck{(= affine), \add{d'anneau $A$},
et $S'$ affine sur $S$ [en remplaçant $S'$ par \ill{} \ill{} \ill{} \ill{}
\ill{}], d'anneau $A'$. On a $A' = \varinjlim A_i$, les $A_i$ étant des
sous-algèbres de type fini de $A'$, et $A'' = A' \otimes_A A' = \varinjlim
A'_i \otimes_A A'$. Comme $H'$ provient d'un $H'_i$ défini sur un $A'_i$,
\uncertain{également} rigide [?], il suffit de \ill{} que les deux images
inverses de $H'_i$ dans $G''_i = G \times_A A''_i$, où $A''_i = A'_i \otimes
A'_i$, sont les mêmes. Comme $A''_i$ est noethérien, il suffit pour cela de}
\marginal{O.K.\ dans le cas de schémas abéliens \ill{} \ill{}, \ill{} \ill{}
\ill{} \uncertain{additifs}}

\ill{} elles coïncident sur la diagonale de $S' \times_S S'$, donc sur toute
composante connexe de $S' \times_S S'$ qui contient la diagonale par
\ill{}. Comme les fibres de $S' \times_S S'$ sont connexes [car celles de $S'$
sont géom.\ connexes] toute composante connexe de $S''$ contient \ill{} \ill{}
\ill{} une fibre, donc
\note{la phrase se poursuit p.~48.}

\page{48}\note{p.~7 de l'auteur.}
rencontre la diagonale, cqfd.

\textbf{Corollaire 3.} Soit $K'$ une extension \emph{primaire} d'un corps
$k$, $G$ un schéma en groupes comm.\ de type fini sur $k$, $H'$ un
sous-schéma \add{en groupes fermé} de $G' = G \otimes_k k'$, \emph{rigide}.
Alors il existe un sous-schéma \add{en groupes} fermé $H$ de $G$, unique,
\struck{rigide} tel que $H' = H \otimes_k k'$.

On admettra qu'il existe une sous-algèbre \struck{$A_1$} de $K'$ \add{de}
\struck{\ill{}} type fini sur $k$, telle que $H'$ provienne d'un sous-schéma
en groupes fermé $H'_1$ de $G_1 = G \otimes_k A_1$, \struck{\ill{}} qui est
\emph{rigide} [c'est le cas en tout cas si \struck{\ill{}} \ill{} que $H'$
\ill{}, sur la clôture algébrique de $K'$, une suite de composition dont les
facteurs sont des $G_m$, Ab, $\mathbb{Z}/\ell\mathbb{Z}$ …].
\marginal{précision à prouver}

\struck{On peut \ill{} \ill{}} Comme le corps $K_1$ de $A_1$ est une extension
primaire de $k$, on voit que $\operatorname{Spec}(A_1)$ est
\add{\uncertain{géométriquement}} \struck{\ill{}} \ill{}. On applique alors le
cor.~2.

\page{49}
\textbf{Cor.~4.} Soit $S' \xrightarrow{f} S$ un morphisme de préschémas
\emph{loc.\ noeth.} On suppose que $f$ est fid.\ plat et de type fini, et
que $S'$ et $S$ sont des spectres de corps \struck{\ill{}} $k'$, $k$
(car $k'/k$ est une extension primaire). \struck{Alors la \ill{}} Soient
\struck{$H$ \ill{} $H \times_S S'$} $H$, $G$ des schémas en groupes
\add{de type fini} commutatifs sur $S$, $H$ étant rigide,
$H' = H \times_S S'$, $G' = G \times_S S'$. Alors l'application canonique
\[
  \operatorname{Hom}_{S\text{-gr}}(H, G) \longrightarrow
  \operatorname{Hom}_{S'\text{-gr}}(H', G')
\]
est bijective. En particulier, le foncteur $H \mapsto H'$ de la catégorie
des schémas \add{de type fini} en groupes commutatifs rigides sur $S$
\struck{dans} \ill{} \ill{} la catégorie \struck{des schémas en groupes
commutatifs} \ill{} \ill{} \ill{} sur $S'$ est pleinement fidèle.

\page{50}\note{p.~8 de l'auteur.}
\marginal{notion : \emph{très rigide}}
Je \ill{} \ill{} \ill{}, dans la définition de la rigidité \add{sur un corps},
d'exiger que tous les sous-schémas fermés \add{\uncertain{rel.}} connexes de
$G$ aient la propriété de \ill{}. Ceci signifie alors que $G$ a une suite de
composition (sur $k$ \emph{alg.\ clos}) \struck{\ill{}} dont les facteurs sont
$\mathbb{Z}/\ell$, $G_m$, Ab, — et \ill{} \add{\uncertain{extensions}} \ill{}
\emph{très rigides}, \ill{} \ill{} \ill{} \ill{}, \ill{} \ill{} bien \ill{}
\ill{} \ill{}, \ill{} \ill{} \ill{} \ill{} \ill{} \ill{} \ill{} \ill{}
\ill{} sur $k$.

\note{la proposition suivante et le début de sa démonstration sont barrés de
plusieurs longs traits obliques.}
\struck{\textbf{Proposition.} $G$ rigide sur $S$ loc.\ noeth., $u \colon G \to
H$ un morphisme de $S$-groupes, \add{\ill{}} où $H$ est un schéma en groupes
de type fini sur $S$. Soit $N = \operatorname{Ker} u$. Alors $N$ est
universellement \ill{} sur $S$.}

\struck{On peut supposer $S = \operatorname{Spec} V$, $V$ anneau} de valuation
discrète.

\textbf{Remarque.} Même si $A$ est un schéma abélien sur un anneau de
valuation discrète \add{\ill{} \ill{} fini}, $B_1$ et $B_2$ deux
sous-schémas \ill{} de $A$ \add{\uncertain{étales}} \struck{\ill{}} sur $V$,
de degré $p$, ayant \ill{} fibres \ill{} \ill{} \ill{}, il est possible que
$B_1 \neq B_2$.

\page{51}
Pour le voir, on prend une courbe elliptique \struck{$B$} $C$ sur $S$, telle
que \ill{} \ill{} $C$ \ill{} \ill{} \ill{} \ill{} de \ill{} \ill{} \ill{}
$C(K)$ \ill{} d'ordre $p$, \ill{} $C(k)$ \ill{} \ill{} $p$, \ill{} \ill{}. \ill{}
\ill{} \ill{} \ill{} elliptique $C'$ sur $S$ \ill{}, \ill{} groupe \ill{} :
$(\mathbb{Z}/p)_S$, \ill{} \ill{} section $\xi$ de $C'$. On prend
$A = C' \times_V C$, on y considère les sections $\alpha = (\xi, e)$,
$\beta = (\xi, \eta)$. Les sous-groupes \add{$\Gamma_1$, $\Gamma_2$}
engendrés sont des sous-groupes de $A$ isomorphes à $(\mathbb{Z}/p)_S$
[\ill{} \ill{} \ill{} \ill{} \ill{} $(\mathbb{Z}/p)_S \to A$, \ill{} \ill{}
\ill{} \ill{} \ill{} \ill{} \ill{} $\mathrm{pr}_1 \colon A \to C'$ \ill{}]. La
\ill{} \ill{} $=$ \struck{\ill{}} fibre en $s$, mais \ill{} \ill{} la fibre
générique $s_1$. Ici, \ill{} \ill{} \ill{}
$A \to A/\Gamma_1 \times A/\Gamma_2$ \struck{\ill{}} \ill{} \ill{} \ill{}
\ill{} sur $\operatorname{Spec} V$, \ill{} \ill{}
\ill{} ${}_s(\mathbb{Z}/p)_{\ker}$ [\ill{} \ill{} \ill{} \ill{} \ill{} \ill{}]
\ill{} \ill{} \ill{} \ill{} \ill{} \ill{} \ill{} \ill{} \ill{} \ill{} de
rigidité.

\page{52}\note{p.~9 de l'auteur.}
\textbf{Proposition.} Soient $k$ un corps, $k'$ une extension
\emph{primaire} de $k$, $G$ un \add{\ill{}} groupe de type fini sur $k$,
$G' = G \otimes_k k'$, $H'$ un sous-groupe fermé de $G'$, tel que
\struck{$H'_0$} \add{$H'^{(0)}_e$} (la composante connexe \ill{} de $e$)
\struck{\ill{}} soit rigide. Alors il existe une \struck{sous-groupe}
extension finie radicielle $k_1$ de $k$, \struck{telle que} et deux
sous-groupes \add{\uncertain{fermés}} \struck{\ill{}} $H_1$ de
$G_1 = G \otimes_k k_1$, tel que $H'_1 = H' \otimes_k k_1$
\struck{\ill{}} \add{\ill{}} $H_1 \otimes_{k_1} k'_1$ (où
$k'_1 = k' \otimes_k k_1$) \ill{} \ill{} \ill{} \ill{}.
\struck{Alors $H'_1 \simeq H_1 \otimes_{k_1} k'_1$.}
\marginal{N.B. Si on veut \ill{} $H'$ \ill{} \ill{} \ill{} \ill{} \ill{}
\ill{} \ill{} extension $K'_1$}

\textbf{Cor.~0.} Si $K'/k$ est \struck{\ill{}} primaire séparable, \ill{}
$K'_1$ \ill{}

\textbf{Corollaire 1.} Il existe une \struck{\ill{}} \struck{\ill{}}
$k \to k_1$ \ill{} \ill{} \ill{} $H$ de $G$ telle que $|H'| =$ image inverse
de $H$.

\textbf{Corollaire 2.} Si $k$ est parfait, il existe un sous-groupe $H$ de
$G$ tel que $H' \supseteq H \otimes_k k'$ et $H' =$ \ill{} \ill{} \ill{}.

\textbf{Corollaire 3.} Soit \struck{$S'$ de type fini sur $S$} \add{$S' \to S$
de type fini surjectif}, \add{\ill{}} $S' \to S$ à fibres rel.\ \ill{}
\ill{} \ill{} irréductibles, \struck{\ill{}} \struck{\ill{} universellement
\ill{} \ill{} \ill{}} $G$ \struck{\ill{}} \add{schéma en groupes commutatif}
de type fini sur $S$, $H'$ un sous-schéma en groupes de
$G' = G \times_S S'$, \add{sur $S'$} tel que $H'$ \ill{} \ill{} sur $S'$
\ill{} \ill{} \ill{} \ill{} fibre \ill{} \ill{} $H'$ aux pts génériques des
fibres de $S'$ \ill{} \add{\uncertain{rigide}}.
\marginal{\struck{hypothèse \ill{} \ill{}}}

\page{53}
Sous ces conditions, il existe une partie \struck{fermée} $H$ de $G$ dont
l'image inverse dans $G'$ soit $H'$. (Dans \ill{} \ill{} morphisme
$S' \to S$ \add{ty.\ surjectif,} \struck{\ill{} \ill{} $S$ \ill{} \ill{}
\ill{}} \ill{} \ill{} \ill{} propre, \ill{} \ill{} \ill{} $S$ loc.\ noeth.)
alors $H$ est une partie fermée de $G$. En \ill{} cas, c'est
« \ill{} \ill{} sous-groupe ».
\marginal{\textbf{Cor.~4.} Sous les conditions \ill{} \ill{} \ill{} cor.~3,
\ill{} \ill{} \ill{} \ill{} \ill{} \ill{}. Alors il existe \ill{} sous-groupe
\ill{} $H$ de $G$ tel que $H' = H \times_S S'$ \ill{}}

Le corollaire 3 résulte facilement du corollaire 2, en \ill{} \ill{} \ill{}
\ill{} \ill{} \ill{} \ill{} $S = \operatorname{Spec}(k)$, et \ill{} \ill{}
\ill{} \ill{} \ill{} corps \struck{\ill{}} \ill{} \ill{} \ill{} \ill{} de $S'$.

\struck{\ill{}}
Démonstration de \struck{la} proposition. \struck{\ill{}} Appliquons le th.\
de rigidité à $H'_0$ ; \ill{} \ill{} \ill{} il existe un sous-groupe $H_0$
de $G$ tel que $H'_0 = H_0 \otimes_k k'$. \ill{} \ill{} \ill{} \ill{} $H_0$,
\ill{} \ill{} \ill{} \ill{} \ill{} \ill{} $H'$ \ill{} un groupe \emph{fini}.
Il \ill{} \ill{} \ill{} \ill{} \ill{} \ill{} \ill{} \ill{} $n$, \ill{} \ill{}
\ill{} ${}_nG' = ({}_nG) \otimes_k k'$. \ill{} \ill{} \ill{} \ill{} \ill{}
\ill{} $G$ est fini sur $k$. \ill{}

\page{54}\note{p.~10 de l'auteur.}
\struck{utilisé dans la}
\note{le lemme qui suit est encadré et barré de trois traits obliques.}
\struck{\textbf{Lemme.} Soit $X$ fini sur un corps $k$, $K'$ une extension
primaire de $k$, $X'$ un sous-schéma \add{(\ill{} fini)} de
$X' = X \otimes_k k'$. Alors \ill{} application \ill{} \ill{}}
\add{et $G/k$ fini}.

Comme $k'/k$ est primaire, l'application $G' \to G$ est bijective (\ill{}
\ill{} \emph{\ill{}}) donc \ill{} \ill{} \ill{} \ill{} $H'$ est l'image
inverse d'un \ill{} \ill{} \ill{} $H$ de $G$. \add{D'ailleurs,}
\struck{\ill{} \ill{} \ill{} \ill{} \ill{}} \ill{} \struck{\ill{}} une
extension radicielle finie convenable $k_1$ de $k$, \struck{\ill{}}
\add{($G \otimes_k k'$ \ill{} \ill{} séparable)} \ill{} $k_1$(\ill{}) des
extensions résiduelles \struck{\ill{}} de $G_1 = G \otimes_k k_1$ \ill{}
séparables. \struck{En} Si donc $H_1$ est l'image inverse de $H$ \ill{}
$G_1$, \ill{} \ill{} \ill{} \ill{} structure induite réduite, $H_1$ est un
sous-schéma \add{en groupes} \emph{fermé} de $G_1$. De plus,
$H_1 \otimes_{k_1} k'_1$ \ill{} $=$ \ill{} \ill{} \ill{} \ill{} que
$H'_1 = H' \otimes_k k'_1$. \ill{} \ill{}.

\page{55}
\textbf{Remarque.} L'hyp.\ sur $H'$ est certainement vérifiée si $G_0$ est
\struck{\ill{}} très rigide. (\uncertain{Cf.}\ ex.\ \ill{} \ill{} \ill{},
\ill{} \ill{} \ill{} \ill{})

\struck{Sous les conditions du cor.~4}, i.e.
\struck{Cor.~5 (\ill{} $H'$ plat sur $S'$, \ill{})}

\textbf{Proposition.} $S'$ de type fini \add{\ill{}} (relativement
\emph{irréductible}) sur $S$, \struck{\ill{}} $G$ schéma en groupes
commutatif de type fini sur $k$, $G' = G \times_S S'$, $H'$ un sous-schéma
en groupes fermé de $G'$ qui est \struck{\ill{}} \add{(\ill{} \ill{} \ill{})}
plat sur $S'$, et dont la composante connexe $H'^{(0)}_\eta$ de la fibre
générique est \emph{rigide} (\ill{} \ill{} \ill{}, et \ill{} que $G$
\add{\uncertain{très}} soit rigide). Alors
\begin{enumerate}
  \item[(i)] il existe un sous-schéma en groupes $H$ de $G$, \ill{} sur $k$,
  tel que $H \otimes_k k'$ \ill{} soit majoré par $H'$ et \ill{} \ill{} \ill{}
  \ill{}. \struck{\ill{} \ill{} \ill{}} \add{\ill{}} \struck{\ill{}} \ill{}
  les conditions \ill{} \ill{} \ill{} \ill{} \ill{} \ill{}
  $u \colon G \to G/H = G_1$ et le morphisme correspondant
  \struck{$G' = G \otimes k' \to$ \ill{}}
  $u' \colon G \otimes_k k' \to G \otimes_k k' / H \otimes_k k' = G'_1$.
  \item[(ii)] \ill{} il
\end{enumerate}
\note{la phrase se poursuit p.~56.}

\page{56}\note{p.~11 de l'auteur.}
existe un unique sous-schéma en groupes \add{fermé} $H'_1$ de $G'_1$, tel
que $H' = u'^{-1}(H'_1)$. \struck{\ill{}, plat} \struck{Et} $H'_1$ est
\add{\ill{}} \struck{fini} et radiciel sur $S'$ \struck{\ill{} \ill{} \ill{}
\ill{} \ill{} \ill{} \ill{} \ill{}} \ill{} \ill{} \ill{} $S'$, il \ill{} \ill{}
\ill{} \ill{} $H'$ \ill{} \ill{} \ill{}. \struck{$H'$ \ill{} \ill{} $H'$ \ill{}
\ill{} \ill{}}
En \ill{} termes, \ill{} \ill{} \ill{} \struck{\ill{}} \ill{} \ill{} \ill{} de
sous-groupes de $G$, \ill{} \ill{} $S'$, dont la \ill{} générique soit une
comp.\ connexe rigide (condition inutile si $G$ est très rigide) est \ill{}
\ill{} \ill{} \ill{} \ill{} \struck{\ill{}} familles plates, \ill{} \ill{}
$S$, de sous-groupes radiciels de \ill{} fixé d'un \ill{} \ill{} $G/H$.
\struck{\ill{} \ill{} \ill{} \ill{}}
\marginal{(iii) \struck{\ill{}} Cela \ill{} \ill{} \ill{} \ill{} \ill{}
\ill{} \ill{} l'espace \ill{} \ill{} \emph{Schéma} \ill{}}
\marginal{($H'_1$ \ill{} \ill{} \ill{}) Si \ill{} \ill{} \emph{Corollaire.}
\ill{} \ill{} \ill{} \ill{} $H'$ \ill{} fibres \ill{} $S$ \ill{}
$H' \simeq H \otimes$ \ill{}}

\emph{Démonstration de la proposition.} \struck{\ill{}} \ill{} \ill{}
\struck{\ill{}} \ill{} cor.~4, \add{(ii)} \ill{} de la proposition plus
générale :

\page{57}
\textbf{\struck{Proposition}} \add{\textbf{Lemme}} $S$ \struck{schéma}
\add{préschéma loc.\ noeth.} \struck{\ill{}} $G$, \add{$G_1$,} préschémas en
groupes \add{de type fini} sur $S$, \struck{\ill{}} \add{\ill{} \ill{}}
plats sur $S$, $u \colon G \to G_1$ un morphisme \emph{surjectif}, (donc
nécessairement plat, \uncertain{surj.}), $N$ son noyau (qui est plat
\add{donc} sur $S$). Alors il y a correspondance biunivoque entre les
sous-schémas fermés $H$ de $G$ qui sont invariants par $N$, et les
sous-schémas fermés \add{$H_1$} de $G_1$, où $H_1$ correspondant à
$H = u^{-1}(H_1)$. Pour que $H$ soit \struck{propre} plat sur $S$, il faut
et il suffit que $H_1$ le soit. Pour que $H$ soit un sous-groupe de $G$, il
faut et il suffit que $H_1$ soit un sous-groupe de $G_1$. Pour que
\add{\uncertain{de plus}} $H_1$ soit radiciel surjectif sur $S$, il faut et
il suffit que $H$ et $N$ aient \add{\ill{}} \ill{} sous-groupes.
\struck{Alors} \ill{} $N \to S$ \add{\ill{} $H \to S$} \add{universellement}
(\ill{}) c'est un morphisme fini radiciel surjectif.
\marginal{\ill{} \ill{} \ill{} \ill{} \ill{} \ill{}}

\uncertain{AQT} par théorie de descente.

\section{Albanese}
\note{titre pris de la p.~58, feuillet par ailleurs blanc portant de sa main
le seul mot « Albanese ».}

\page{59}
\emph{Schéma d'Albanese.}

Soit $X$ \struck{projectif et} \add{fid.} plat \add{qu.-proj.} sur $S$,
\struck{$\mathcal{O}_S \xrightarrow{\sim} f_*(\mathcal{O}_X)$ universel}.

\struck{On dit} \struck{Pour une espèce} \uncertain{Données équivalentes} :
\begin{enumerate}
  \item[(i)] $P$ fibré principal homogène [\ill{} \ill{} \ill{} \ill{} de
  \uncertain{qu.-proj.}] sur $S$, \ill{} groupe $G$, et un \ill{}
  morphisme $\varphi \colon X \to P$ ;
  \item[(i bis)] Un $S$-groupe $G$, et une donnée de descente \ill{} sur
  $G_X = G \times_S X$ par $X \to S$, et la structure de fibré pr.\ homog.\
  sur $G_X$ \add{\ill{}} \ill{} \ill{} $G_X$ ;
  \item[(ii)] Un $S$-groupe $G$, et un schéma
  $\varphi \colon X \times_S X \to G$ satisfaisant la condition
  \add{\uncertain{symétrique}}
  $\varphi(x,y)\varphi(y,z) = \varphi(x,z)$ [i.e.\ \ill{} \ill{} \ill{}
  $\varphi(x,x) = e$], et une condition d'effectivité ; [automatiquement
  satisfaite si $G$ est affine$/S$, ou si \ill{} $X \to$ \uncertain{qu.\ pr.}\
  sur $S$, \ill{} \ill{} \ill{} \ill{} \ill{} \ill{}, et $X \to S$ admet loc.\
  des quasi-sections finies, p.~ex.\ $X \to S$ projectif].
\end{enumerate}
\note{en marge gauche, en face de (i bis), un carré :}

\begin{tikzcd}
P \arrow[d] & P_X \arrow[l] \arrow[d] \\
S & X \arrow[l] \arrow[ul, dashed, "\varphi"']
\end{tikzcd}

\note{la flèche $\varphi$ est figurée par un point et une lettre au centre du
carré ; sa direction est une lecture. En marge, en face de (ii) :}
$X \leftleftarrows X \times_S X$ \ill{} $X \times_S X \times_S X$.

Supposons enfin que \struck{\ill{}} $X$ soit \struck{\ill{}}
\uncertain{projectif} sur $S$ et tel que
$\mathcal{O}_S \xrightarrow{\sim} f_*(\mathcal{O}_X)$ universel ; \ill{}
\ill{} \ill{} \ill{} \ill{}
\begin{enumerate}
  \item[(iii)] Un $S$-groupe $G$, et un schéma
  $\varphi \colon X \times_S X$ \ill{} satisfaisant $\varphi(x.x) = e$, et
  \ill{} \ill{} \ill{}. \ill{} \ill{}$/S$ et $X/e$ \ill{} \ill{} \ill{}
  \ill{} $X$ projectif \ill{}.
\end{enumerate}
\marginal{définition \ill{} $G$ \ill{} \ill{} \ill{} $(P, u)$ \ill{} \ill{}
\ill{} $G$ \ill{} \ill{} \ill{} \ill{} \ill{}
$\operatorname{Hom}_{S'}(X_{S'}, G_{S'})$ \ill{} \ill{}
$\operatorname{Pic}(X_{S'}, G_{S'})$ \ill{} \ill{} … et \ill{} \ill{}
associé …}

Pour simplifier, nous supposons que la suite $f \colon X \to S$
\struck{projectif}, et \ill{} $\mathcal{O}_S \simeq f_*(\mathcal{O}_X)$
universel ! \ill{} \ill{} \add{\ill{}} \ill{} \ill{} \ill{} \ill{} $G$
\add{\uncertain{proj.}} sur $S$, \ill{} \ill{} \ill{} fibré principal
homogène \ill{} $P$, \ill{} \ill{} \ill{} \ill{} \ill{} \ill{} \ill{} \ill{}
$\operatorname{Hom}_S(X, P)$ comme foncteur \ill{} \ill{} \ill{} $P$,
\add{\ill{} \ill{} $T_X(P)$}. \ill{} \ill{} \ill{}, ce foncteur est
\emph{pro-représentable} \ill{} la catégorie des $P$, \ill{} \ill{}
\emph{\ill{}}. \ill{} \ill{} \ill{} \ill{}

\page{60}
\ill{} existent, et que \ill{} \ill{} \ill{} \ill{} représentable : \ill{}
\ill{} \ill{} $T_X$ \ill{} \ill{} : \ill{}, \ill{} \ill{} \ill{} \ill{}
\ill{} \ill{} \ill{} couple $(P, \varphi \in T_X(P))$ \ill{} \ill{} \ill{}
\struck{\ill{}} donné \ill{} \ill{} \ill{} \ill{} « couple universel ».

Interprétation, en \ill{} \ill{} \ill{} \ill{} \ill{} \ill{}
$\varphi \colon X \times_S X \to G$ satisfaisant $\varphi(x,x) = e$, alors
\struck{\ill{}} les \ill{} \ill{} \ill{} groupes \ill{} $G_\alpha$ \ill{}
\ill{} tels que $\varphi$ se factorise par $G_\alpha$ \ill{} forment un
\ill{} filtrant décroissant, et \ill{} \ill{} \ill{} il y a un \ill{} \ill{}
minimal, soit $G'$. Le couple $(G', \varphi')$ correspondant est \ill{}
\ill{} cherché.

\struck{Montrons que \ill{} \ill{} \ill{} \ill{} qui représente $T_X$ \ill{}
\ill{} \ill{} \ill{} $(G_\alpha, P_\alpha)$ \ill{} \ill{}.} \ill{} \ill{}.
Mais \ill{} \ill{} \ill{} \ill{} \ill{} \ill{} \ill{} \ill{} $G$ à être un
schéma abélien $A$, \ill{} $A = \operatorname{Pic}^\circ_{B/S}$, où
$B = \hat{A} = \operatorname{Pic}^\circ_{A/S}$. En effet, \ill{} \ill{}
\ill{} $S' = X$, \ill{} \ill{} \ill{} \struck{soit $T_X(B)$ \ill{} \ill{}
\ill{}}
\begin{align*}
  T_X(A) &\simeq \operatorname{Hom}_S(X \times_S X, A)
  \simeq \operatorname{Hom}_{S'\text{-pointés}}(X_{S'}, A_{S'}) \\
  &= \operatorname{Hom}_{S'\text{-pointés}}(X_{S'},
     \operatorname{Pic}^\circ_{B_{S'}/S'}) \\
  &\simeq \mathcal{B}_{S'}(X_{S'}, B_{S'})
  \simeq \operatorname{Hom}_{S'\text{-groupes}}(B_{S'},
     \operatorname{Pic}_{X_{S'}/S'}) \\
  &\simeq \operatorname{Hom}_{S\text{-groupes}}(B,
     \operatorname{Pic}_{X/S})
\end{align*}
\note{sous $\mathcal{B}_{S'}(X_{S'}, B_{S'})$, accolade : « données de corr.\
divisorielles » ; sous le dernier $\simeq$ : « d'après la platitude et la
rigidité … ». La première ligne se prolonge en marge droite par
$= \operatorname{Hom}_{S'\text{-gr}}(B_{S'}, \operatorname{Pic}_{\ill{}})$,
illisible.}

Encadré :
\[
  T_X(A) \simeq \operatorname{Hom}_{S\text{-gr}}(\hat{A},
  \operatorname{Pic}_{X/S})
\]

\ill{} est un résultat du type de « \uncertain{dualité} de \ill{} ».
\add{On suppose que $\operatorname{Pic}_{X/S}$} contient \ill{} \ill{},
\ill{} \ill{} ($B_{X/S}$ \ill{} \ill{} \ill{} \ill{}
$\operatorname{Pic}^\circ_{X/S}$), \ill{} \ill{} « \ill{} d'\ill{} \ill{} »
\note{la page s'arrête ici ; la suite est hors du lot.}

\end{document}
